\documentclass[11pt]{article}

\usepackage[margin=25mm]{geometry}

\usepackage{amsmath,amssymb,fontspec,unicode-math}

\IfFontExistsTF{DejaVu Serif}{\setmainfont{DejaVu Serif}}{\setmainfont{Latin Modern Roman}}

\IfFontExistsTF{Latin Modern Math}{\setmathfont{Latin Modern Math}}{\setmathfont{DejaVu Math TeX Gyre}}

\usepackage{xurl,hyperref}

\setlength{\parindent}{0pt}\setlength{\parskip}{6pt}

\begin{document}

\title{Equilibrium Policy Generalization: A Reinforcement Learning Framework for Cross-Graph Zero-Shot Generalization in Pursuit-Evasion Games}

\author{ELVIA preliminary AI analysis}\date{21 October 2026}\maketitle

Preliminary assessment: Minor.\par

Independent scientific review remains pending. Author adoption has not been requested.

The public dynamic program’s cost depends on predecessor degrees; equilibrium alone does not ensure finite capture. Evaluation-distribution and sampled-minimum differences need separate interpretation.

\section*{Review scope}

Latest-version, claim-focused review of the cited author-paper passages and applicable current author-release paths. The stated counterexamples and scope controls were checked analytically; no author program, old verifier, training run or benchmark was executed.

\section*{Problem 1: Queue uniqueness does not bound predecessor-query work}

Source: Section 3.2.1, p. 6; Algorithm 1, p. 5; official DP\_implementation (2 pursuers).cpp.

Assessment: supported; impact: minor

\subsection*{Source claim (paraphrase)}

The implementation’s runtime argument counts each state being enqueued once but does not retain its predecessor-neighborhood costs.

\subsection*{Analysis}

On a complete graph with two pursuers and stay actions, all \(n^3\) states are terminal at initialization. Each popped state still inspects \(n\) evader and \(n^2\) pursuer predecessors, yielding order \(n^6\) candidate checks. Unsuccessful scans cost work even when they add no queue entry.

\subsection*{Finding and scope}

The released degree-independent near-linear cost is not established. Bounded-degree behavior or an improved indexed algorithm is not excluded.

\subsection*{Proposed correction}

State dependence on agent count and maximum degree, or supply and analyze the promised implicit neighborhood structure including unsuccessful queries.

An explicit local correction is described; author adoption and a corresponding published erratum are not established.

\section*{Problem 2: Equilibrium existence does not imply finite capture everywhere}

Source: Section 3.2.1, p. 6, sentence following complexity paragraph.

Assessment: supported; impact: minor

\subsection*{Source claim (paraphrase)}

The theoretical discussion infers guaranteed finite capture from pure-strategy equilibrium without restricting to a capture-winning region.

\subsection*{Analysis}

Use three isolated vertices with stay actions, distinct starts for two pursuers and one evader. There is only one possible policy pair, hence a pure equilibrium, but capture never occurs. The model does not impose connectedness or a universally winning region.

\subsection*{Finding and scope}

The finite-capture implication is too broad. A value identity with \(\gamma^{\infty}=0\) can remain valid.

\subsection*{Proposed correction}

Restrict finite-capture claims to the finite capture-winning region and state the infinite-rank convention separately.

An explicit local correction is described; author adoption and a corresponding published erratum are not established.

\section*{Problem 3: Imported and new curves use different test distributions}

Source: Current arXiv v2 Section 4.3, PDF p.10, Figure 4; Appendix D.2, differing baseline initial-game filters

Assessment: supported; impact: minor

\subsection*{Source claim (paraphrase)}

Appendix D.2 discloses that the original baseline filtering could not be reproduced and substitutes distance-filtered games for the new evaluation.

\subsection*{Analysis}

The imported and new Figure 4 curves therefore average different initial-game distributions. Differences in average difficulty can affect scores even without a policy change. The disclosed choice is not hidden leakage, and no ranking reversal is established.

\subsection*{Finding and scope}

The figure juxtaposes protocols; it does not isolate a controlled policy effect.

\subsection*{Proposed correction}

Evaluate both methods on shared graph edges, starts, exits, horizon and reward/discount conventions, or state the different-protocol scope explicitly.

An explicit local correction is described; author adoption and a corresponding published erratum are not established.

\section*{Problem 4: A minimum of sampled returns is not a minimum of expected returns}

Source: Official supplement multi-exit/5vs1-pursuer/test\_worker.py run\_episode\_shortest and stochastic driver.

Assessment: supported; impact: minor

\subsection*{Source claim (paraphrase)}

The stochastic worker minimizes one realized return per exit, whereas worst-case expected utility is \(\min_e\mathbb E[R_e]\).

\subsection*{Analysis}

In the stated connected five-pursuer, eight-exit example, each exit is captured with probability \(1/8\). Return \(+1\) for capture and \(-1\) otherwise gives worst expected utility \(-3/4\). Independent single draws followed by a minimum instead have expectation \(2(1/8)^8-1\approx-0.99999988\). In general \(\mathbb E[\min_eR_e]\leq\min_e\mathbb E[R_e]\). Deterministic-return evaluation is unaffected.

\subsection*{Finding and scope}

The stochastic release statistic targets a different quantity. No actual published model score is recalculated.

\subsection*{Proposed correction}

Estimate the conditional mean for each exit using repeated returns, then minimize those estimates and quantify finite-sample selection uncertainty.

An explicit local correction is described; author adoption and a corresponding published erratum are not established.

\section*{Limits}

Preliminary AI assessment; independent human scientific adjudication remains pending.\par

An analytic counterexample or a release discrepancy does not by itself remeasure the paper’s reported experiments.\par

\section*{Source identity}

arXiv:2511.00811v2 (current author version; updated 2025-12-12T06:37:26Z)

\url{https://arxiv.org/pdf/2511.00811v2}

Reviewed PDF SHA-256: \nolinkurl{0a9e0ff848e3f972d6a93d71610edbbb7a665978f64404dfd603a1ddde046083}

The selected author/publisher source is the same work as the original case, matched by exact title, author identity and official version channel. The exact selected PDF SHA-256 was recomputed in the latest-version pass for Issue \#27. The PDF is a retained author/publisher copy; its original retrieval time is stated only when supported by matching acquisition metadata. File modification time was not substituted.

\end{document}
